% File: numodel-manual.tex
% Standalone manual for the numodel package.
% Source for the documentation that used to live in numodel.dtx.

\documentclass{ltxdoc}
\usepackage{numodel}
\usepackage[left=3.5cm, right=2cm, top=2cm, bottom=2cm,
            marginparwidth=3.5cm, marginparsep=0.3cm]{geometry}
\usepackage{xcolor}
\usepackage{booktabs}
\definecolor{numorange}{rgb}{0.90,0.45,0.00}
\EnableCrossrefs
\CodelineIndex

% Fonts.  fontspec + unicode-math require LuaLaTeX (already the
% engine for this bundle).  unicode-math is loaded AFTER numodel
% (and thus after amsmath) so it integrates with the existing math
% setup rather than fighting it.
\usepackage{fontspec}
\setmainfont{Arial}
\usepackage{unicode-math}
\setmathfont{Lete Sans Math}
\setmonofont{Fira Mono}

% Make ltxdoc's |...| inline verbatim breakable: long option strings
% and key paths (e.g. |/pgf/number format/use comma|) overflow the
% margin otherwise.  fvextra's \DefineShortVerb with
% breaklines+breakanywhere replaces the doc.sty default.
%
% NB: ltxdoc.cls binds | via \AtBeginDocument{\MakeShortVerb{\|}}.
% That hook runs at \begin{document}, after the preamble.  So the
% rebinding must happen INSIDE a later \AtBeginDocument hook --
% otherwise ltxdoc's binding overwrites fvextra's.
\usepackage{fvextra}
\AtBeginDocument{%
  \MakeShortVerb{\|}% match ltxdoc's first, so DeleteShortVerb succeeds
  \DeleteShortVerb{\|}%
  \DefineShortVerb[breaklines,breakanywhere]{\|}%
}

% tcolorbox with the listings library: lets us write a single
% model-example body that is both displayed verbatim AND executed
% (instead of duplicating the code into a verbatim block + a
% redeclaration).  See \begin{modelexample}...\end{modelexample}
% below.
\usepackage{tcolorbox}
\tcbuselibrary{listings,skins,breakable}

\lstdefinestyle{numodelcode}{
  basicstyle=\small\ttfamily,
  breaklines=true,
  columns=fullflexible,
  keepspaces=true,
  showstringspaces=false,
}

\NewTCBListing{modelexample}{}{%
  enhanced,
  breakable,
  listing style=numodelcode,
  colback=black!3,
  colframe=black!50,
  boxrule=0.4pt,
  arc=2pt,
  before skip=10pt,
  after skip=10pt,
}

% Match tick-label backgrounds to the modelexample colback so plots
% emitted by \diagrammodel blend into the listing's gray frame
% instead of the numodel-plot default (fill=white).
\pgfplotsset{numodel/axis/.append style={%
  xticklabel style={fill=black!3},%
  yticklabel style={fill=black!3}%
}}

\begin{document}

\CheckSum{0}

\changes{v0.1}{2026/04/24}{Initial version, extracted from internal
  project sources.}
\changes{v0.2.0}{2026/05/16}{l3build workflow; bundle structure
  with numodel-plot; diagram-style key; units key; localised
  column titles; decimal-separator key; factor-aware flow
  detection.  See CHANGELOG.md for details.}
\changes{v0.3.0}{2026/05/17}{Version-sync release with
  \textsf{numodel-plot}~v0.3.0; no functional changes to
  \textsf{numodel} itself.}
\changes{v0.4.0}{2026/05/19}{Expression reference table; display
  translation for sqrt/exp/ln/sin/cos/tan/asin/acos; \cs{textmodel}
  migrated to \textsf{tabularray}'s \texttt{longtblr} with
  localised continuation markers; new \texttt{tblrenv} key for
  \cs{textmodel}, \cs{numodelsetup} and package option.  See
  CHANGELOG.md for details.}
\changes{v0.5.0}{2026/05/23}{Multi-series \cs{diagrammodel}
  (comma-separated y-variables, unit-aware filtering, colour-blind
  safe palette); shared inflow/outflow Forrester valves with curved
  branches; improved flow-variable heuristic (aux/stock wins over
  constant in inflow terms; top-level parens distributed before
  classification).  See CHANGELOG.md for details.}
\changes{v0.7.0}{2026/05/28}{Multi-prefix \cs{diagrammodel}:
  \texttt{prefix=} accepts a comma-separated list of prefixes; the
  same y-variables are plotted across every prefix, with the prefix
  appended in parentheses after each legend entry.  Fixed display of
  multi-character numeric superscripts in \cs{textmodel}
  (\texttt{T\textasciicircum 1.2}, \texttt{T\textasciicircum 12},
  \texttt{T\textasciicircum -3}) — now wrapped in braces so the
  whole exponent typesets as a superscript.  See CHANGELOG.md for
  details.}

\GetFileInfo{numodel.sty}

\DoNotIndex{\newcommand,\newenvironment,\def,\edef,\let,\global,
  \RequirePackage,\ProvidesPackage,\NeedsTeXFormat,\endinput,
  \ExplSyntaxOn,\ExplSyntaxOff,\begin,\end,\relax,\undefined,\cs_new,
  \cs_set,\cs_gset,\tl_new,\tl_set,\int_new,\seq_new,\prop_new}

\title{The \textsf{numodel} package\thanks{This document corresponds
  to \textsf{numodel}~\fileversion, dated \today.}}
\author{Paul Zuurbier \\ \texttt{mail@paulzuurbier.nl}}
\date{\today}
\maketitle

\begin{abstract}
A LuaLaTeX package for writing and rendering numerical models
(Euler-integrated dynamical systems) directly inside LaTeX
documents, aimed at physics teaching material.  It provides a
text-model pipeline (\texttt{\textbackslash mvar},
\texttt{\textbackslash mrule}, \texttt{\textbackslash computemodel}),
Forrester stock-and-flow diagrams, and optional plots of the
computed time series via the sibling package \textsf{numodel-plot}.
\end{abstract}

\tableofcontents

\section{Introduction}

\textsf{numodel} lets an author write a dynamical system (stocks,
flows, helper variables, rules, and a stop condition) as a sequence
of LaTeX macros and renders three complementary views of that model
directly in the document:
\begin{itemize}
\item a \emph{text model} \textemdash{} a typeset rule table with
  initial values (|\textmodel|);
\item a \emph{graphic model} \textemdash{} a Forrester
  stock-and-flow diagram with auto-layout (|\graphicmodel|);
\item a \emph{diagram} \textemdash{} a numerical Euler simulation
  plus PGFPlot of any pair of variables (|\computemodel| followed
  by |\diagrammodel|; the plot is rendered through the sibling
  package \textsf{numodel-plot}).
\end{itemize}

All three views are produced from a single set of declarations so
the textbook description, the conceptual stock-and-flow diagram,
and the numerical result of the same model are guaranteed to stay
in sync.  Variables and rules live in namespaces (\emph{prefixes})
so a document can contain multiple independent models;
|\newmodelprefix{P}| starts a fresh one.  The simulation engine
runs in Lua (through \textsf{luacode}) for $\mathcal{O}(1)$ appends
and cheap min/max tracking; the rendering layer is pure
\textsf{expl3}.

\section{First example: a free-falling ball}

A ball dropped from $h_0 = 100\,\mathrm{m}$ under constant
gravitational acceleration.  The complete model:

\begin{quote}
\begin{verbatim}
\usepackage[syntax=EN]{numodel}

\newmodelprefix{ball}
\mvar{T}{t}{0}{\s}{2}{system}
\mvar{Dt}{dt}{0.1}{\s}{2}{system}
\mvar{V}{v}{0}{\m\per\s}{2}{stock}
\mvar{Y}{y}{100}{\m}{3}{stock}
\mvar{G}{g}{-9.81}{\m\per\s\squared}{3}{aux}

\mrule{V}{\ballV + \ballG * \ballDt}
\mrule{Y}{\ballY + \ballV * \ballDt}
\mrule{T}{\ballT + \ballDt}
\mstop{\ballY <= 0}
\end{verbatim}
\end{quote}

After the declarations above, three render commands produce the
three views shown below, each from the \emph{same} model.

\subsection{\texttt{\textbackslash textmodel} \textemdash{} rule table}

The verbatim source rendered by |\textmodel|:
\newmodelprefix{ball}
\mvar{T}{t}{0}{\s}{2}{system}
\mvar{Dt}{dt}{0.1}{\s}{2}{system}
\mvar{V}{v}{0}{\m\per\s}{2}{stock}
\mvar{Y}{y}{100}{\m}{3}{stock}
\mvar{G}{g}{-9.81}{\m\per\s\squared}{3}{aux}
\mrule{V}{\ballV + \ballG * \ballDt}
\mrule{Y}{\ballY + \ballV * \ballDt}
\mrule{T}{\ballT + \ballDt}
\mstop{\ballY <= 0}

\begin{center}
\textmodel
\end{center}

Each |\mvar| with a non-empty start value contributes a row in the
\emph{initial values} column; each |\mrule| contributes a row in the
\emph{model} column.  Symbols come from the second |\mvar|
argument (the display text), values are formatted through
\textsf{siunitx}.  |<=| is rendered as $\leqslant$.

\subsection{\texttt{\textbackslash graphicmodel} \textemdash{}
Forrester diagram}

|\graphicmodel| draws the same model as a stock-and-flow diagram.
Stocks (type |stock|) are rectangles; auxiliaries (|aux|) and
constants (|constant|) are circles, with two short horizontal
dashes flanking the constant ring.  In the default |tight|
diagram-style, an aux or constant that is the direct inflow or
outflow of a stock is absorbed into the valve label and not drawn
as a separate node; |diagram-style=forrester| or |edu| keep both
(Section~\ref{sec:cfg}).

Flows are inferred per additive term on the right-hand side of
each stock rule.  After dropping the self-carry term, every |+T|
is a candidate inflow and every |-T| a candidate outflow of the
target stock; the first non-|system|, non-|Dt| variable appearing
inside~|T| becomes the flow variable (|aux| and |stock| beat
|constant| when several candidates coexist).  When the same flow
variable surfaces as an outflow term of stock~|A| and a matching
inflow term of stock~|B|, the pair is promoted to a
\emph{between-flow}: one valve drains |A| into |B| instead of two
separate valves.  If the flow variable of an inflow term is itself
a stock, the source stock takes the valve's place -- with a
matching outflow term on its own rule (a conserved-quantity
between-flow) or, when no such match exists, with a cloud-fed
phantom valve plus a causal link to the source stock:

\begin{center}
\graphicmodel
\end{center}

Layout is automatic from the rule graph: stocks (with their
inflow/outflow valves) sit on the bottom row, auxiliaries on the
middle row, and constants on the top row, each filled
left-to-right in declaration order.  The order of |\mvar| calls
therefore drives the horizontal reading order of the diagram --
reordering, inserting or holding back a |\mvar| nudges the
visible arrangement without touching any grid coordinates, which
is the primary lever a document author has over the look of the
diagram.  Manual placement remains available through |gridx|/|gridy|
keys on the |\mvar| call (see Section~\ref{sec:api}).  Wide
diagrams can be capped to a maximum number of grid columns with
|\numodelsetup{gridmaxx=N}|: when a row reaches |N|, the auto
layout shifts the affected items up one row and continues filling.

\subsection{\texttt{\textbackslash computemodel} +
\texttt{\textbackslash diagrammodel} \textemdash{} numerical plot}

|\computemodel| iterates the rules forward in time using
Euler integration with step size~|\ballDt|, stopping when |\mstop|'s
condition becomes true (or when the |maxiter| safety limit is
reached, see Section~\ref{sec:cfg}).  |\diagrammodel{xvar}{yvar}{label}|
then plots one variable against another:

\begin{modelexample}
\computemodel
\diagrammodel{T}{Y}{ball-fall}
\end{modelexample}

Axis ranges, tick lattice, and labels are computed automatically
from the simulated min/max of each variable; the plot inherits the
\textsf{numodel-plot} style (see that package's documentation for
configuration).

\section{Configuration}\label{sec:cfg}

\DescribeMacro{\numodelsetup}
Runtime configuration:
\begin{quote}
\begin{verbatim}
\numodelsetup{syntax=NL, maxiter=50000}
\end{verbatim}
\end{quote}
The same keys can also be passed as package options:
|\usepackage[syntax=NL]{numodel}|.  Recognised keys:

\begin{description}
\item[\texttt{syntax}] Language tag for the rule-table rendering.
  Built-in values:
  \begin{description}
  \item[\texttt{EN}] (default) XMILE-style ALL-CAPS keywords:
    |IF|/|THEN|/|ELSE|, |AND|, |OR|, |ABS|, |SIGN|, \ldots
  \item[\texttt{NL}] Dutch CoachTaal keywords:
    |Als|/|Dan|/|Anders|, |EN|, |OF|, |Abs|, |Teken|, \ldots
  \item[\texttt{FPEVAL}] (aliases |fpeval|, |l3fp|) shows each rule
    body verbatim in l3fp (|\fpeval|) syntax, exactly as typed in
    |\mrule| -- lowercase function names, literal |*|, |^|, |&&|,
    |?:| -- typeset monospace.  Variable macros still resolve to
    their display names; the ternary is \emph{not} expanded into
    |IF|/|THEN|/|ELSE| and the stop row reads |stop:|~\meta{cond}.
    Decimal mark defaults to |point|.
  \end{description}
  Each language tag |X| corresponds
  to a file |numodel-X.def| located via |kpse| when the package
  processes the key.  The package ships with |numodel-EN.def| and
  |numodel-NL.def|; drop your own |numodel-FR.def| (or any other
  tag) in |TEXMFHOME/tex/latex/numodel/| and select it with
  |\usepackage[syntax=FR]{numodel}| -- no package rebuild needed.
  The setting affects display only; the expression syntax in
  |\mrule| bodies is always |\fp_eval|-compatible.
\item[\texttt{maxiter}] Safety limit on the number of |\computemodel|
  iterations (default 20\,000).  When reached, the simulation
  aborts with a warning naming the unmet stop condition.
\item[\texttt{graphscalex}] Horizontal grid spacing in centimetres
  for |\graphicmodel|'s Forrester layout (default 2).  Larger
  values spread the diagram out horizontally.
\item[\texttt{graphscaley}] Vertical grid spacing in centimetres
  for |\graphicmodel|'s Forrester layout (default 2).  Larger
  values spread the diagram out vertically.

  |graphscalex|/|graphscaley| control only the empty space
  between diagram elements: the nodes themselves (stocks,
  valves, auxiliaries, constants) have fixed dimensions
  expressed in |em| and are not individually configurable.
  Because |em| is relative to the current font size, the way
  to grow or shrink the nodes uniformly is to change the font
  size around the |\graphicmodel| call -- e.g.\ wrap it in
  |{\small \graphicmodel}| or |{\large \graphicmodel}|.  Pick
  spacing with |graphscalex|/|graphscaley|, pick node size
  with the surrounding font.
\item[\texttt{gridmaxx}] Maximum number of grid columns the auto
  layout may fill on any one row before wrapping (integer, default
  |0| = no limit).  When the limit is reached, items already placed
  on the affected row (and everything above it for the stocks row,
  everything but stocks for the aux row, only the constants for the
  constants row) shift up by one row to free space, and placement
  continues from column~0.  Manually positioned variables
  (|\mvar[gridx=...,gridy=...]|) are kept where they are.  When
  wrapping is active the default centring of the aux row and the
  right-aligning of the stocks row are disabled, so the diagram fills
  left-to-right, bottom-to-top.
\item[\texttt{diagram-style}] Rendering style for the case where a
  helper or constant is the direct inflow/outflow of a stock.
  Three values:
  \begin{description}
  \item[\texttt{tight}] (default) the valve takes the helper's/constant's
    label; the helper/constant itself is not drawn as a separate node.
    Compact and most LaTeX-native.
  \item[\texttt{forrester}] Forrester/Sterman convention: the valve is
    drawn without a label and the helper/constant remains as a separate
    node connected to the valve by a causal arrow.
  \item[\texttt{edu}] Didactic dual form: the valve carries the label
    \emph{and} the helper/constant is drawn as a separate node with a
    causal arrow to the valve.  Visually busy but pedagogically explicit.
  \end{description}
\item[\texttt{flowarrow-style}] Visual style of the flow pipe.
  |hollow| renders the classic Forrester double-line pipe with an
  open arrow head; |filled| renders a thick solid arrow.  The
  default tracks |diagram-style|: |forrester| picks |hollow|, the
  other styles pick |filled|.  An explicit value overrides this
  coupling.
\item[\texttt{valve-style}] Visual style of the valve node.
  |valve| draws the bow-tie/butterfly icon (Forrester); |circle|
  draws an empty circle on the flow pipe; |edu| draws a labelled
  circle (the flow variable's display text inside).  The default
  tracks |diagram-style|: |forrester| picks |valve|, the other
  styles pick |edu|.
\item[\texttt{flowarrow-cloud-tip}] Whether the open end of an
  inflow or outflow pipe is anchored to a cloud node, signalling
  the model boundary.  Default tracks |diagram-style|: |forrester|
  picks |true|, the other styles pick |false|.  May be set
  globally via |\numodelsetup|, per-render via |\graphicmodel|, or
  per-stock via |\mvar[flowarrow-cloud-tip=...]|.  The most
  specific source wins.
\item[\texttt{units}] Whether the \emph{initial values} cells in
  |\textmodel| display the SI unit alongside the value (|\qty|) or
  only the numeric value (|\num|).  Boolean, default |true|.  May
  also be supplied to |\textmodel[units=false]| as a per-table
  override; the global setting is restored after rendering.
\item[\texttt{tblrenv}] Which \textsf{tabularray} environment wraps
  the |\textmodel| table.  Three values: |longtblr| (default,
  page-breakable); |tblr| (inline, no page breaks); |talltblr|
  (inline, no page breaks, supports |\caption|/notes).  Pick |tblr|
  or |talltblr| when |\textmodel| sits inside an enclosing
  environment that suppresses page breaks (|subfigure|, |minipage|,
  \dots); the default |longtblr| would otherwise emit spurious
  ``(Continued)'' / \emph{Continued on next page} markers in that
  setting.  May also be supplied per render as
  |\textmodel[tblrenv=...]|; the global setting is restored
  afterwards.
\item[\texttt{decimal-separator}] Decimal mark used by every number
  that |numodel| renders: the \emph{initial values} column of
  |\textmodel|, and the tick labels of |\diagrammodel|.  Two values:
  \begin{description}
  \item[\texttt{comma}] use a comma (\textsf{siunitx}
    |output-decimal-marker={,}|, \textsf{pgfplots}
    |/pgf/number format/use comma|).
  \item[\texttt{point}] use a full stop (\textsf{siunitx}
    |output-decimal-marker={.}|, \textsf{pgfplots}
    |/pgf/number format/use period|).
  \end{description}
  The default tracks |syntax|: |NL| picks |comma|, |EN| and
  |FPEVAL| pick |point|.  Other language files can publish a default by
  defining |\__numodel_kw_<LANG>_dsep_default:| (expanding to
  |point| or |comma|); when the macro is absent the default is
  |point|.  An explicit |decimal-separator| key locks that
  choice and overrides any future |syntax| change.  The override is
  scoped: |numodel| applies it only inside its own renderers
  (\textsf{siunitx}'s state is restored on group exit), so a
  document-wide |\sisetup| is not perturbed.
\end{description}

\subsection{Diagram styles in practice}

The same model rendered under each of the three |diagram-style|
values.  The model is the simplest case that distinguishes the
styles: one stock $N$ with a constant inflow~$R$:

\begin{modelexample}
\newmodelprefix{flux}
\mvar{T}{t}{0}{\s}{2}{system}
\mvar{Dt}{dt}{1}{\s}{2}{system}
\mvar{N}{n}{0}{}{0}{stock}
\mvar{R}{r}{5}{\per\s}{2}{constant}
\mrule{N}{\fluxN + \fluxR * \fluxDt}
\mrule{T}{\fluxT + \fluxDt}
\mstop{\fluxT >= 5}
\begin{tabular}{@{}ccc@{}}
\graphicmodel[diagram-style=tight] &
\graphicmodel[diagram-style=forrester] &
\graphicmodel[diagram-style=edu] \\[2pt]
\texttt{tight} & \texttt{forrester} & \texttt{edu}
\end{tabular}
\end{modelexample}

\begin{itemize}
\item \texttt{tight} collapses the constant $R$ into the valve
  label, producing the most compact diagram.
\item \texttt{forrester} keeps the canonical System-Dynamics
  convention: unlabelled bow-tie valve, the constant remains a
  separate node, the link from $R$ to the valve is a thin causal
  arrow.
\item \texttt{edu} is a didactic dual: the valve carries the label
  \emph{and} the constant remains as a separate node with a causal
  arrow.  Less compact but pedagogically explicit -- useful when
  first introducing the stock/flow vocabulary.
\end{itemize}

\section{Public API}\label{sec:api}

\subsection{Variables and rules}

\DescribeMacro{\mvar}
Declares a model variable.  Signature:
\begin{quote}
|\mvar[<keys>]{<Name>}{<text>}{<start>}{<unit>}{<sig>}{<type>}|
\end{quote}
where \meta{Name} is a short alphabetic identifier (the
prefix-qualified accessor becomes |\<prefix><Name>|, along with a
family of suffixed accessor macros |\<prefix><Name><suffix>|
documented in Section~\ref{sec:accessors}), \meta{text} is the
math-mode display symbol used in the rule table and diagram
(e.g.\ |F_{res}|), \meta{start} is the initial value (a number,
or empty for helpers computed by a rule, or any \textsf{expl3}
\texttt{fp}-evaluable expression involving previously defined
model variables), \meta{unit} is a bare \textsf{siunitx} unit
macro sequence (e.g.\ |\m\per\s\squared|), \meta{sig} is the
number of significant figures used by |\<prefix><Name>num|/|qty|,
and \meta{type} is one of |stock|, |aux|, |constant|, or
|system|.  Each English type also accepts a Dutch alias
(|voorraad|, |hulp|, |constante|, |systeem|) for backwards
compatibility with existing teaching material.  See
Section~\ref{sec:types}.

\textbf{Naming caveat.}  All accessor macros of a given prefix
share one flat TeX namespace: |\mvar{X}| generates |\<prefix>X|
and every |\<prefix>X<suffix>| listed in
Section~\ref{sec:accessors}.  Pick \meta{Name}s so that no name
equals another name followed by such a suffix, otherwise the two
definitions silently overwrite each other.  The canonical trap is
declaring both |T| and |Tmax|: the |max|-suffix accessor of |T|
(the extremum |\<prefix>Tmax|) collides with the no-suffix
accessor of |Tmax|, so |\<prefix>Tmax| inside a rule expression
ends up meaning whichever was declared last.  The same hazard
applies to prefix choices -- avoid prefixes where one is a
concatenation of another with a variable name (e.g.\ |ball| with
a variable |Y| versus a prefix |ballY|).  An |\mvar| name of
|steps| would similarly collide with the per-prefix iteration
accessor |\<prefix>steps| set by |\computemodel|.

Optional \meta{keys}:
\begin{description}
\item[\texttt{prefix}] Override the current prefix for this single
  call.
\item[\texttt{gridx}, \texttt{gridy}] Manual placement in the
  |\graphicmodel| grid; integers, $-1$ leaves the slot to
  auto-layout (default).
\item[\texttt{alias}] Math-mode token list that replaces the
  entire \emph{initial values} cell.
\item[\texttt{aliasleft}, \texttt{aliasright}] Replace just the
  left symbol or right value half of the cell.
\end{description}

\DescribeMacro{\mrule}
Adds a rule of the form $\meta{LHS} \leftarrow \meta{expr}$.
Signature:
\begin{quote}
|\mrule*[<keys>]{<LHS>}{<expr>}|
\end{quote}
Both forms add the rule to \emph{both} the rule table
(|\textmodel|) and the simulation (|\computemodel|); execution is
identical.  The star only changes the typeset layout when
\meta{expr} is a ternary |cond ? a : b|.  Without the star the
ternary is rendered on a single table row,
\begin{quote}
\begin{verbatim}
IF cond THEN lhs = a ELSE lhs = b ENDIF
\end{verbatim}
\end{quote}
which is compact but wide.  With the star (|\mrule*|) the same
ternary is broken across rows,
\begin{quote}
\begin{verbatim}
IF cond THEN
    lhs = a
ELSE
    lhs = b 
ENDIF
\end{verbatim}
\end{quote}
keeping the table column narrow so a |\graphicmodel| can sit
alongside it, and making the source itself easier to read.  For
non-ternary expressions the star has no effect.
\meta{expr} may use the full |\fp_eval| expression grammar.
See Table~\ref{tab:expr} (Section~\ref{sec:expr}) for a complete
overview of supported operators and functions with their XMILE and
CoachTaal equivalents.

\DescribeMacro{\mruletext}
Inserts a free-text row in the rule table without registering a
rule with the simulator.  Signature: |\mruletext[<keys>]{<text>}|.
Useful for inserting comments or section dividers in long
rule tables.

\DescribeMacro{\mstop}
Sets the simulation stop condition.  Signature:
|\mstop[<keys>]{<expr>}|.  The simulation halts at the first step
where \meta{expr} evaluates true.  Exactly one |\mstop| per model
prefix is required before |\computemodel|.  Without one,
|\computemodel| issues a warning.

\subsection{Expression reference}\label{sec:expr}

Table~\ref{tab:expr} lists all expression constructs for |\mrule|
bodies.  The \textbf{XMILE} column shows XMILE\,v1.0 syntax,
\textbf{l3fp} shows the |\fp_eval|-compatible form used inside
|\mrule|, and \textbf{CoachTaal} shows the Coach\,7 equivalent
(function names from the standard CoachTaal math reference;
note that argument separators in CoachTaal are semicolons).
\textcolor{numorange}{Orange} entries are not yet fully supported
by \textsf{numodel} (the expression computes correctly, but the
rendered keyword in the rule table is not yet translated).
\textit{Italic} cells contain a derived equivalent rather than a
native keyword.

\begin{longtblr}[
  caption = {Expression reference: XMILE, \texttt{\textbackslash fp\_eval},
             and CoachTaal},
  label   = {tab:expr},
]{
  colspec = {@{} l l l @{}},
  rowhead = 1,
  rowsep  = 1pt,
}
  \textbf{XMILE} & \textbf{l3fp (\texttt{\textbackslash fpeval})}
                 & \textbf{CoachTaal} \\
  \hline

  \SetCell[c=3]{l} \itshape Arithmetic operators \\
  \hline
  \texttt{+} & \texttt{+} & \texttt{+} \\
  \texttt{-} & \texttt{-} & \texttt{-} \\
  \texttt{*} & \texttt{*} & \texttt{*} \\
  \texttt{/} & \texttt{/} & \texttt{/} \\
  \texttt{\textasciicircum} &
    \texttt{\textasciicircum} &
    \texttt{\textasciicircum} \\
  \textcolor{numorange}{\texttt{MOD(x,y)}}
    & \textcolor{numorange}{\textit{x - trunc(x/y)*y}}
    & \textcolor{numorange}{\textit{x - Entier(x/y)*y}} \\

  \SetCell[c=3]{l} \itshape Comparison operators \\
  \hline
  \texttt{<}  & \texttt{<}  & \texttt{<}  \\
  \texttt{<=} & \texttt{<=} & \texttt{<=} \\
  \texttt{>}  & \texttt{>}  & \texttt{>}  \\
  \texttt{>=} & \texttt{>=} & \texttt{>=} \\
  \texttt{=}  & \texttt{=}  & \texttt{=}  \\
  \texttt{<>} & \texttt{!=} & \texttt{<>} \\

  \SetCell[c=3]{l} \itshape Boolean operators \\
  \hline
  \texttt{AND} & \texttt{\&\&} & \texttt{EN}   \\
  \texttt{OR}  & \texttt{\textbar\textbar} & \texttt{OF}   \\
  \texttt{NOT} & \texttt{!}    & \texttt{NIET} \\

  \SetCell[c=3]{l} \itshape Control flow \\
  \hline
  \texttt{IF c THEN a ELSE b}
    & \texttt{c ? a : b}
    & \texttt{Als c Dan a Anders b EindAls} \\

  \SetCell[c=3]{l} \itshape General math functions \\
  \hline
  \texttt{ABS(x)}  & \texttt{abs(x)}  & \texttt{Abs(x)}  \\
  \texttt{SIGN(x)} & \texttt{sign(x)} & \texttt{Teken(x)} \\
  \texttt{SQRT(x)} & \texttt{sqrt(x)} & \texttt{Sqrt(x)} \\
  \textcolor{numorange}{\texttt{INT(x)}}
    & \textcolor{numorange}{\texttt{trunc(x)}}
    & \textcolor{numorange}{\textit{Entier(Abs(x))*Teken(x)}} \\
  \textcolor{numorange}{\textit{INT(x+0.5)}}
    & \textcolor{numorange}{\texttt{round(x)}}
    & \textcolor{numorange}{\texttt{Round(x)}} \\
  \textit{INT(x)} & \texttt{floor(x)} & \texttt{Entier(x)} \\
  \textit{-INT(-x)} & \texttt{ceil(x)} & \textit{-Entier(-x)} \\
  \texttt{MIN(x,y)} & \texttt{min(x,y)} & \texttt{Min(x;y)} \\
  \texttt{MAX(x,y)} & \texttt{max(x,y)} & \texttt{Max(x;y)} \\
  \textcolor{numorange}{\textit{---}}
    & \textcolor{numorange}{\texttt{fact(x)}}
    & \textcolor{numorange}{\texttt{Fac(x)}} \\
  \textcolor{numorange}{\textit{INT(LOG10(ABS(x)))}}
    & \textcolor{numorange}{\texttt{logb(x)}}
    & \textcolor{numorange}{\textit{Entier(Log(Abs(x)))}} \\

  \SetCell[c=3]{l} \itshape Exponential and logarithmic \\
  \hline
  \texttt{EXP(x)} & \texttt{exp(x)} & \texttt{Exp(x)} \\
  \texttt{LN(x)}  & \texttt{ln(x)}  & \texttt{Ln(x)}  \\
  \textcolor{numorange}{\texttt{LOG10(x)}}
    & \textcolor{numorange}{\textit{ln(x)/ln(10)}}
    & \textcolor{numorange}{\texttt{Log(x)}} \\

  \SetCell[c=3]{l} \itshape Trigonometry (radians) \\
  \hline
  \texttt{SIN(x)}    & \texttt{sin(x)}  & \texttt{Sin(x)}    \\
  \texttt{COS(x)}    & \texttt{cos(x)}  & \texttt{Cos(x)}    \\
  \texttt{TAN(x)}    & \texttt{tan(x)}  & \texttt{Tan(x)}    \\
  \texttt{ARCSIN(x)} & \texttt{asin(x)} & \texttt{Arcsin(x)} \\
  \texttt{ARCCOS(x)} & \texttt{acos(x)} & \texttt{Arccos(x)} \\
  \texttt{ARCTAN(x)} & \texttt{atan(x)} & \texttt{Arctan(x)} \\
  \textcolor{numorange}{\texttt{ARCTAN2(y,x)}}
    & \textcolor{numorange}{\texttt{atan(y,x)}}
    & \textcolor{numorange}{\textit{Arctan(y/x)}} \\
  \textit{1/TAN(x)} & \texttt{cot(x)} & \textit{1/Tan(x)} \\
  \textit{1/SIN(x)} & \texttt{csc(x)} & \textit{1/Sin(x)} \\
  \textit{1/COS(x)} & \texttt{sec(x)} & \textit{1/Cos(x)} \\
  \textit{ARCSIN(1/x)} & \texttt{acsc(x)} & \textit{Arcsin(1/x)} \\
  \textit{ARCCOS(1/x)} & \texttt{asec(x)} & \textit{Arccos(1/x)} \\
  \textit{ARCTAN(1/x)} & \texttt{acot(x)} & \textit{Arctan(1/x)} \\
  \textcolor{numorange}{\textit{ARCTAN2(x,y)}}
    & \textcolor{numorange}{\texttt{acot(y,x)}}
    & \textcolor{numorange}{\textit{Arctan(x/y)}} \\

  \SetCell[c=3]{l} \itshape Trigonometry (degrees) \\
  \hline
  \textit{SIN(PI/180*x)} & \texttt{sind(x)} & \textit{Sin(Pi/180*x)} \\
  \textit{COS(PI/180*x)} & \texttt{cosd(x)} & \textit{Cos(Pi/180*x)} \\
  \textit{TAN(PI/180*x)} & \texttt{tand(x)} & \textit{Tan(Pi/180*x)} \\
  \textit{180/PI*ARCSIN(x)} & \texttt{asind(x)} & \textit{180/Pi*Arcsin(x)} \\
  \textit{180/PI*ARCCOS(x)} & \texttt{acosd(x)} & \textit{180/Pi*Arccos(x)} \\
  \textit{180/PI*ARCTAN(x)} & \texttt{atand(x)} & \textit{180/Pi*Arctan(x)} \\
  \textcolor{numorange}{\textit{180/PI*ARCTAN2(y,x)}}
    & \textcolor{numorange}{\texttt{atand(y,x)}}
    & \textcolor{numorange}{\textit{180/Pi*Arctan(y/x)}} \\
  \textit{1/TAN(PI/180*x)} & \texttt{cotd(x)} & \textit{1/Tan(Pi/180*x)} \\
  \textit{1/SIN(PI/180*x)} & \texttt{cscd(x)} & \textit{1/Sin(Pi/180*x)} \\
  \textit{1/COS(PI/180*x)} & \texttt{secd(x)} & \textit{1/Cos(Pi/180*x)} \\
  \textit{180/PI*ARCSIN(1/x)} & \texttt{acscd(x)} & \textit{180/Pi*Arcsin(1/x)} \\
  \textit{180/PI*ARCCOS(1/x)} & \texttt{asecd(x)} & \textit{180/Pi*Arccos(1/x)} \\
  \textit{180/PI*ARCTAN(1/x)} & \texttt{acotd(x)} & \textit{180/Pi*Arctan(1/x)} \\
  \textcolor{numorange}{\textit{180/PI*ARCTAN2(x,y)}}
    & \textcolor{numorange}{\texttt{acotd(y,x)}}
    & \textcolor{numorange}{\textit{180/Pi*Arctan(x/y)}} \\

  \SetCell[c=3]{l} \itshape Constants \\
  \hline
  \texttt{PI} & \texttt{pi} & \texttt{Pi} \\
  \texttt{e}  & \textit{exp(1)}  & \textit{Exp(1)}  \\
  \textcolor{numorange}{\texttt{INF}}
    & \textcolor{numorange}{\texttt{inf}}
    & \textcolor{numorange}{\textit{---}} \\
  \textcolor{numorange}{\texttt{NAN}}
    & \textcolor{numorange}{\texttt{nan}}
    & \textcolor{numorange}{\textit{---}} \\
  \textcolor{numorange}{\textit{PI/180}}
    & \textcolor{numorange}{\texttt{deg}}
    & \textcolor{numorange}{\textit{Pi/180}} \\
  \textcolor{numorange}{\textit{1}}
    & \textcolor{numorange}{\texttt{true}}
    & \textcolor{numorange}{\texttt{Aan}} \\
  \textcolor{numorange}{\textit{0}}
    & \textcolor{numorange}{\texttt{false}}
    & \textcolor{numorange}{\texttt{Uit}} \\

  \SetCell[c=3]{l} \itshape Simulation-specific functions (XMILE only) \\
  \hline
  \textcolor{numorange}{\texttt{TIME}}
    & \textcolor{numorange}{\textit{user variable (\texttt{\textbackslash mvar})}}
    & \textcolor{numorange}{\textit{user variable}} \\
  \textcolor{numorange}{\texttt{DT}}
    & \textcolor{numorange}{\textit{user variable (\texttt{\textbackslash mvar})}}
    & \textcolor{numorange}{\textit{user variable}} \\
  \textcolor{numorange}{\texttt{STEP(h,t\textsubscript{0})}}
    & \textcolor{numorange}{\textit{T >= t0 ? h : 0}}
    & \textcolor{numorange}{\textit{Als T >= t0 Dan h Anders 0 EindAls}} \\
  \textcolor{numorange}{\texttt{RAMP(s,t\textsubscript{0})}}
    & \textcolor{numorange}{\textit{T > t0 ? s*(T-t0) : 0}}
    & \textcolor{numorange}{\textit{Als T > t0 Dan s*(T-t0) Anders 0 EindAls}} \\
  \textcolor{numorange}{\texttt{DELAY(x,dt)}}
    & \textcolor{numorange}{\textit{not supported}}
    & \textcolor{numorange}{\textit{not supported}} \\
  \textcolor{numorange}{\texttt{SMOOTH(x,t)}}
    & \textcolor{numorange}{\textit{not supported}}
    & \textcolor{numorange}{\textit{not supported}} \\
  \textcolor{numorange}{\texttt{INIT(x)}}
    & \textcolor{numorange}{\textit{not supported}}
    & \textcolor{numorange}{\textit{not supported}} \\
  \textcolor{numorange}{\texttt{PREVIOUS(x)}}
    & \textcolor{numorange}{\textit{not supported}}
    & \textcolor{numorange}{\textit{not supported}} \\
  \textcolor{numorange}{\textit{RANDOM(0,1)}}
    & \textcolor{numorange}{\texttt{rand()}}
    & \textcolor{numorange}{\texttt{Rand}} \\
  \textcolor{numorange}{\texttt{RANDOM(lo,hi)}}
    & \textcolor{numorange}{\textit{lo + (hi-lo)*rand()}}
    & \textcolor{numorange}{\textit{lo + (hi-lo)*Rand}} \\
  \textcolor{numorange}{\textit{not supported}}
    & \textcolor{numorange}{\texttt{randint(n)}}
    & \textcolor{numorange}{\textit{not supported}} \\
  \textcolor{numorange}{\textit{not supported}}
    & \textcolor{numorange}{\texttt{randint(m,n)}}
    & \textcolor{numorange}{\textit{not supported}} \\
  \hline
\end{longtblr}

\subsection{Render commands}

\DescribeMacro{\textmodel}
Renders the rule-and-startvalue table.  Optional |[<keys>]| accepts
|prefix=<name>| (render a non-current model), |units=true| or
|units=false|, and |tblrenv=tblr|longtblr|talltblr|.  Each per-call
key overrides the global |\numodelsetup| setting for this single
render only; the global state is restored afterwards.

|tblrenv| selects which |tabularray| environment wraps the table:
|longtblr| (default) breaks across pages, |tblr| renders inline
without page breaks, |talltblr| renders inline but with caption and
note support.  Use |tblr| or |talltblr| when |\textmodel| sits inside
an enclosing environment that suppresses page breaks (|subfigure|,
|minipage|, \dots); the default |longtblr| would otherwise emit
spurious continuation markers there.

\textbf{Row spacing.}  |numodel| loads |tabularray| and sets
|\SetTblrInner{rowsep=0pt}| globally so the rule listing renders
compactly.  This applies to \emph{every} |tabularray| table in the
document; if you want the default spacing back in your own tables,
issue |\SetTblrInner{rowsep=2pt}| (or whatever value you prefer)
somewhere in your document.

\DescribeMacro{\graphicmodel}
Renders the Forrester stock-and-flow diagram.  Variables of type
|stock| become rectangles, |constant| become circles, |aux|
become identifier nodes; flow arrows connect stocks to constant or
helper sources/sinks based on which variables appear in
the right-hand side of stock-updating rules.

Optional |[<keys>]| accepts |prefix=<name>| (render a non-current
model) and |diagram-style=tight|forrester|edu|, which overrides the
global |\numodelsetup| setting for this single render only.  The
global state is restored afterwards, so multiple |\graphicmodel|
calls can each pick their own style without re-issuing
|\numodelsetup|.

\DescribeMacro{\computemodel}
Runs the Euler simulation.  Each iteration step is executed in
LaTeX itself: \textsf{expl3}'s |\fp_eval| (the engine behind the
LaTeX2e |\fpeval|) evaluates the stop condition and then every
|\mrule| body in declaration order, writing each new value into
the accessor macro |\<prefix><Name>|.  The per-step values are
appended to Lua tables on the side -- storing the time series as
growing TeX token lists would cost $O(N)$ per append and $O(N^2)$
overall, whereas the Lua-table append is $O(1)$, so a 20\,000-step
run stays linear in total work.  Lua also keeps the running min
and max so no second pass over the series is needed.  After
|\computemodel| returns, |\<prefix><Name>min| /
|\<prefix><Name>max| hold the extrema, |\<prefix><Name>| holds
the final-step value, |\<prefix>steps| expands to the number of
recorded samples (i.e.\ the iteration count $N$), and the full
time-series can be retrieved with |\mcoords| / |\mstep|.

\DescribeMacro{\diagrammodel}
Convenience wrapper that produces a complete |figure| with caption
and label.  Signature:
\begin{quote}
|\diagrammodel[<keys>]{<xvar>}{<yvar>|\meta{,\dots}|}[<extra>]{<label>}|
\end{quote}
Reads min/max and display text from |\<prefix><xvar>| etc.,
delegates to |\drawplot| from \textsf{numodel-plot}, and emits
|\caption{$<prefix><yvar><text>(<prefix><xvar><text>)$-diagram}\label{fig:<label>}|.  The optional
\meta{extra} argument is appended to the |axis| body, useful for
additional |\addplot| lines (annotations, theoretical curves).
Must be called after |\computemodel|.

The \meta{yvar} argument accepts a comma-separated list of
variables: every entry whose |unitraw| matches the first one's is
drawn into the same diagram as discrete model points (no connecting
lines), with the y-axis range scaled to the joint min/max of the
kept series.  Entries with a non-matching unit are dropped and a
warning is issued.  The seven colours cycle through Okabe \& Ito's
colour-blind safe palette (yellow omitted), ordered so consecutive
series differ in luminance -- which keeps them distinguishable on a
greyscale printout as well.  The legend lists each kept variable's
display text.

The |prefix=| key likewise accepts a comma-separated list, so the
same y-variables are plotted across several models in one diagram --
useful for direct visual comparison between alternative model
variants of the same physical setup.  Series order is prefix-major
(all y-variables of the first prefix come before any of the second),
the colour cycle continues across the cross product, and axis
ranges reduce over the union of all (prefix, yvar) pairs.  The unit
filter is applied against the first prefix's first y-variable, so
the same short names are assumed to share unit and display text
across prefixes.  When more than one prefix is supplied the prefix
is appended in parentheses after each y-variable's display text in
the legend, e.g.\ |$U_3$ (ptc)| and |$U_3$ (pw)|.  Single-prefix
calls render exactly as before.

\begin{quote}
|\diagrammodel[prefix={ptc,pw}]{T}{U_3,U_par}{compare-models}|
\end{quote}

Example -- two quantities in one diagram.  The free-fall ball from
the introduction is extended with gravitational potential energy
$E_p = m\,g\,h$ (written as $-m\,G\,Y$ because $G = -9{,}81$ already
encodes the downward direction) and kinetic energy
$E_k = \tfrac{1}{2}\,m\,v^2$.  Both share the unit |\joule|, so the
multi-series filter keeps both and scales the y-axis to their
combined range.  At $t=0$ all energy sits in $E_p$; as the ball
falls, $E_p$ drops and $E_k$ grows.  The two |aux| variables
receive a start-value expression so that step~0 is recorded with
the physically meaningful initial energies rather than empty
values.
\begin{modelexample}
\newmodelprefix{nrg}
\mvar{T}{t}{0}{\s}{2}{system}
\mvar{Dt}{dt}{0.1}{\s}{2}{system}
\mvar{V}{v}{0}{\m\per\s}{2}{stock}
\mvar{Y}{y}{100}{\m}{3}{stock}
\mvar{m}{m}{0.5}{\kg}{3}{constant}
\mvar{G}{g}{-9.81}{\m\per\s\squared}{3}{constant}
\mvar{Ep}{E_p}{-\nrgm * \nrgG * \nrgY}{\joule}{3}{aux}
\mvar{Ek}{E_k}{0.5 * \nrgm * \nrgV * \nrgV}{\joule}{3}{aux}
\mrule{V}{\nrgV + \nrgG * \nrgDt}
\mrule{Y}{\nrgY + \nrgV * \nrgDt}
\mrule{Ep}{-\nrgm * \nrgG * \nrgY}
\mrule{Ek}{0.5 * \nrgm * \nrgV * \nrgV}
\mrule{T}{\nrgT + \nrgDt}
\mstop{\nrgY <= 0}
\computemodel
\diagrammodel{T}{Ep,Ek}{ball-energy}
\end{modelexample}

\subsection{Series accessors}

\DescribeMacro{\mcoords}
Returns a comma-separated PGFPlots coordinate list of the
simulated series for two variables.  Fully expandable.  Signature:
|\mcoords{<xvar>}{<yvar>}| (current prefix);
|\mcoordsp{<prefix>}{<xvar>}{<yvar>}| (explicit prefix).  Both
forms are usable inside |\addplot coordinates{ ... }|.

\DescribeMacro{\mstep}
Returns the value of one variable at a chosen iteration.  Fully
expandable.  Signature: |\mstep{<Name>}{<i>}|;
|\mstepp{<prefix>}{<Name>}{<i>}|.  The current prefix is prepended
automatically by |\mstep|.  Step indexing is 0-based: step~0 is
the initial-values row, step~$N{-}1$ is the final recorded step
after |\computemodel|.  Negative indices count from the end
(Python-style), so |\mstep{Y}{-1}| is the last recorded $y$ and
|\mstep{Y}{-2}| is the penultimate one.  Raises a |numodel| error
when \meta{i} falls outside the recorded range
(|0..\<prefix>steps - 1| for non-negative indices, or
|-\<prefix>steps..-1| for the Python-style negative form), which
surfaces typo'd indices and "ran $N$ steps but indexed step $N$"
off-by-one mistakes rather than expanding to nothing and propagating
empty arguments into the surrounding plot.

Example -- a red secant line through the last two simulated
$(t,y)$ points of the free-fall ball (the model from the first
example), extrapolated across the whole $t$-domain and labelled in
the legend.  The slope is the rise-over-run of the last two steps,
the intercept is the final $y$-value.  Inside the optional
|[<extra>]| argument |\mstep| expands as a literal number into
PGFPlots' math parser, so each |\mstep| is evaluated only once
(when the expression is constructed) rather than per sample.
Applied to the |ball| model:

\switchmodelprefix{ball}
\begin{modelexample}
\diagrammodel{T}{Y}[%
  \addplot[red, very thick, domain=\ballTmin:\ballTmax]
    {\mstep{Y}{-1}
      + (\mstep{Y}{-1} - \mstep{Y}{-2})
      / (\mstep{T}{-1} - \mstep{T}{-2})
      * (x - \mstep{T}{-1})};
  \addlegendentry{secant endpoints}
]{ball-fall-secant}
\end{modelexample}

\subsection{Namespace management}

\DescribeMacro{\newmodelprefix}
Creates a new model namespace and switches to it.  Signature:
|\newmodelprefix{<name>}|.  Subsequent |\mvar| / |\mrule| / |\mstop|
calls bind to this prefix.  All generated accessors are prefixed
(so |\mvar{Y}{y}{...}{...}{...}{...}| under prefix |ball| produces
|\ballY|, |\ballYtext|, etc.).

\DescribeMacro{\switchmodelprefix}
Switches to a previously created prefix.  Useful when a document
defines several models early and renders them later out of order.
Signature: |\switchmodelprefix{<name>}|.

\DescribeMacro{\NumodelForEachVar}
Iterates a token list over every registered variable across every
known prefix.  Signature: |\NumodelForEachVar{<code with #1>}|;
inside the body, |#1| is the full prefixed name (e.g.\ |ballY|).
Used by external tools (e.g.\ a worksheet system) that need to
expand every model accessor.

\subsection{Shared valves}\label{sec:shared-valve}

When the same auxiliary variable appears as the inflow of more
than one stock, |\graphicmodel| draws a single valve next to the
first such stock and threads a curved branch (|to[bend left=30]|
-- the same bend as the curved causal arrows) from that valve over
the primary stock into each additional one.  The auto-layout
leaves the involved stocks side by side on the same row so the
branch stays compact.  A minimal example:
\begin{quote}
\begin{verbatim}
\mvar{R}{r}{5}{\m\cubed\per\s}{2}{aux}
\mvar{Va}{V_1}{0}{\m\cubed}{3}{stock}
\mvar{Vb}{V_2}{0}{\m\cubed}{3}{stock}
\mrule{Va}{\Va + \R * \Dt}
\mrule{Vb}{\Vb + \R * \Dt}
\end{verbatim}
\end{quote}
renders one |R| valve, a straight pipe |R| $\to$ |V_a|, and a
curved branch |R| $\to$ |V_b| arcing over |V_a|.  (Note: TeX
control words do not accept digits, so the macro names |Va|/|Vb|
are used internally while the display texts |V_1|/|V_2| appear
in the rendered diagram and table.)

Shared \emph{outflows} are handled symmetrically.  When one
variable drains more than one stock, the valve is placed to the
right of the last-declared source stock, and every earlier
source attaches with a curved branch arcing over the intervening
stocks into the shared valve.  This matters in particular for
physical models where a single Stefan--Boltzmann constant or
friction coefficient drains several reservoirs into the same
sink.

Two flow-classification refinements support these layouts:
\begin{itemize}
\item The flow-variable heuristic prefers \texttt{aux} and
  \texttt{stock} variables over \texttt{constant}s when both
  appear in the same additive term, so an inflow valve carries
  the meaningful rate variable instead of an earlier-declared
  scaling constant.
\item Top-level parenthesised inflow terms of the form
  |(A - B)*C| are distributed into |+A*C, -B*C| before
  classification, so |A| surfaces as the inflow valve and
  |B| (typically a radiative or dissipative loss involving a
  constant) surfaces as the outflow valve.
\end{itemize}

\section{Variable types}\label{sec:types}

The sixth |\mvar| argument (\meta{type}) tags a variable with a
role.  The type drives both |\textmodel| layout (whether the row
appears in \emph{initial values}) and |\graphicmodel| node shape.
Each type has a canonical English name and a Dutch alias; the two
are interchangeable.

\begin{description}
\item[\texttt{stock}] A stock that
  accumulates over time.  Drawn as a rectangle in the Forrester
  diagram; its rule must be of the form
  $\text{stock} \leftarrow \text{stock} + \dots$ so that the
  integrator can detect inflows and outflows.  Receives an
  \emph{initial values} row.
\item[\texttt{constant}] A constant
  parameter (mass, gravitational acceleration, spring stiffness,
  $\ldots$).  Drawn as a circle.  Receives an \emph{initial values}
  row.
\item[\texttt{aux}] An auxiliary variable
  computed from other variables on each step.  Drawn as a plain
  identifier node, no rectangle.  No \emph{initial values} row
  (start value is normally left empty).
\item[\texttt{system}] System-level
  bookkeeping (time |T|, step size |Dt|, terminal time,
  $\ldots$).  Drawn separately, no flow arrows.  Receives an
  \emph{initial values} row.
\end{description}

\section{Generated accessors}\label{sec:accessors}

Each |\mvar| call generates a family of accessor macros named
|\<prefix><Name><suffix>|.  The full set:

\begin{description}
\item[\textit{(no suffix)}] Current numeric value (post-rule-update
  if inside a step, post-final-step after |\computemodel|).
\item[\texttt{text}] The display symbol passed as the second
  argument of |\mvar|.
\item[\texttt{unit}] |\unit{...}| applied to the unit argument
  (\textsf{siunitx}-formatted).
\item[\texttt{unitraw}] The unit argument verbatim, without the
  |\unit{}| wrapper, suitable for use as a building block (e.g.\
  |\xlabelunit| in \textsf{numodel-plot}).
\item[\texttt{num}] The current value formatted via \textsf{siunitx}'s
  |\num{}| with the variable's significant figures.  Used by
  |\textmodel| in the \emph{initial values} column when |units=false|.
\item[\texttt{qty}] As |num|, but including the unit
  (|\qty{value}{unit}|).  Used by |\textmodel| in the
  \emph{initial values} column when |units=true| (the default).
\item[\texttt{pre}] As |qty|, but with engineering prefix mode
  (e.g.\ |1500 W| renders as |1{,}5\,kW|).
\item[\texttt{sign}] Significant-figure count (raw integer).
\item[\texttt{type}] Variable type (raw string).
\item[\texttt{min}, \texttt{max}] Extrema over the simulated series.
  Empty before |\computemodel|.
\item[\texttt{gridx}, \texttt{gridy}] Manual placement coordinates
  in the Forrester grid; $-1$ if left to auto-layout.
\item[\texttt{alias}, \texttt{aliasleft}, \texttt{aliasright}]
  Override tokens for the \emph{initial values} cell layout.
  |aliasright| replaces just the value part of the cell (keeping
  the symbol and equals sign); |aliasleft| replaces just the
  symbol; |alias| replaces the entire cell.
\end{description}

In addition, |\computemodel| writes one per-\emph{prefix} accessor
(no variable name in the middle):
\begin{description}
\item[\texttt{\textbackslash <prefix>steps}] The number of recorded
  samples after the most recent |\computemodel| call for this
  prefix, i.e.\ the iteration count $N$.  Empty (undefined) before
  the first |\computemodel|.  Picking |steps| as an |\mvar| name
  would collide with this accessor -- see the naming caveat in
  Section~\ref{sec:api}.
\end{description}

Example -- the alias keys exist primarily as a worksheet hook:
the same model can be rendered as a fully-specified reference
\emph{and} as a fill-in-the-blanks exercise.  |aliasright=\cdots|
blanks the value while keeping the symbol visible (asking
\emph{what value} belongs there), |alias={?}| blanks the entire
cell (asking the student to supply both the symbol and the value
from physical knowledge).  Crucially, |\computemodel| keeps using
the original numeric start value, so the diagram or computed
answer for the rest of the model remains correct -- only the
rendered table changes.  The row for $y$ now reads \(y = \cdots\)
(asking the student to identify the initial height from a graph
or photograph) and the row for $g$ shows just \(?\) (asking the
student to supply the gravitational acceleration from memory).
The rule column is unaffected -- |\mrule| has its own
|alias|/|aliasright| keys for blanking rule bodies if needed.

\begin{modelexample}
\newmodelprefix{quiz}
\mvar{T}{t}{0}{\s}{2}{system}
\mvar{Dt}{dt}{0.5}{\s}{2}{system}
\mvar{V}{v}{0}{\m\per\s}{2}{stock}
\mvar[aliasright=\cdots]{Y}{y}{50}{\m}{3}{stock}
\mvar[alias={?}]{G}{g}{-9.81}{\m\per\s\squared}{3}{constant}
\mrule{V}{\quizV + \quizG * \quizDt}
\mrule{Y}{\quizY + \quizV * \quizDt}
\mrule{T}{\quizT + \quizDt}
\mstop{\quizY <= 0}
\textmodel
\end{modelexample}

\section{Multiple models in one document}

Each |\newmodelprefix| starts a fresh namespace.  This lets a
document contain several unrelated models without name clashes:
\begin{quote}
\begin{verbatim}
\newmodelprefix{ball}
\mvar{Y}{y}{100}{\m}{3}{stock}
% ... ball model ...

\newmodelprefix{spring}
\mvar{X}{x}{0.1}{\m}{3}{stock}
% ... spring model ...

% Render the ball model:
\switchmodelprefix{ball}\textmodel\computemodel\diagrammodel{T}{Y}{fall}

% Render the spring model:
\switchmodelprefix{spring}\textmodel\computemodel\diagrammodel{T}{X}{spring}
\end{verbatim}
\end{quote}
Each prefix carries an independent rule list, stop condition, and
recorded series.

\section{Requirements}

\textsf{numodel} requires LuaLaTeX (the engine, for the Lua
runtime) and TeX~Live~2022 or later.  Mandatory dependencies:
\textsf{expl3}, \textsf{xparse}, \textsf{l3keys2e}, \textsf{amsmath},
\textsf{amssymb}, \textsf{xfrac}, \textsf{tikz}, \textsf{luacode},
\textsf{siunitx}, \textsf{float}, and the sibling package
\textsf{numodel-plot}
(which itself pulls in \textsf{pgfplots}).  The companion Lua
module |numodel.lua| must be installed alongside the |.sty| in a
directory searched by |kpse|.

\end{document}
